\begin{pastexam}[label={px:ntu-108-im-it-02}]{臺大資管所 108 學年度 問答題第 2 題}
（問答；20 分）

Double hashing is a popular collision resolution scheme that employs two hash functions to
determine the location of an item in a hash table. One hash function computes the starting
location of an item and the other one specifies the size of probe step.
\begin{enumerate}[label=(\alph*)]
\item (10 points) Explain the benefit of double hashing when resolving collisions.
\item (5 points) What is the problem if the hash table size and the probe step size are not
relatively prime?
\item (5 points) Suggest a solution to the problem.
\end{enumerate}
\end{pastexam}

\begin{pastsoln}
記號：表的大小 $m$，第 $i$ 次探測（$i = 0, 1, 2, \ldots$）的位置是
\[
h(k, i) = (h_1(k) + i \cdot h_2(k)) \bmod m ,
\]
$h_1(k)$ 是起點，$h_2(k)$ 是步長（資結書第~12.3~節）。

\textbf{(a) 好處。} 和另外兩種開放定址法比：
\begin{itemize}
\item \textbf{線性探測}（步長固定 $1$）有\textbf{主要群聚}（primary clustering）：已經佔用的
格子連成一長串，落在這一串裡任何位置的新鍵都要走到串尾，串又因此變得更長，越長越容易變長。
\item \textbf{二次探測}（$h_1(k) + c_1 i + c_2 i^2$）打散了長串，但兩個鍵只要起點 $h_1$ 相同，
探測序列就完全一樣，這叫\textbf{次要群聚}（secondary clustering）。
\end{itemize}
雙重雜湊的步長也由鍵決定。 兩個鍵就算 $h_1$ 相同，$h_2$ 通常不同，第二步就分開了，所以主要
和次要群聚都沒有。 例如 $m = 13$、$h_1(k) = k \bmod 13$、$h_2(k) = 1 + (k \bmod 7)$：$k = 14$
和 $k = 27$ 的起點都是 $1$，步長分別是 $1$ 和 $7$，探測序列是 $1, 2, 3, \ldots$ 和
$1, 8, 2, 9, \ldots$，不會擠在同一條路上。 線性探測和二次探測的探測序列由起點決定，只有 $m$
種；雙重雜湊由 $(h_1(k), h_2(k))$ 兩個值決定，可以有大約 $m^2$ 種，行為接近「每個鍵的探測順序
都是隨機排列」的理想情形。 在這個理想模型下，負載因子 $\alpha$ 時失敗搜尋的期望探測次數不超過
$1/(1 - \alpha)$，成功搜尋不超過 $\frac{1}{\alpha}\ln\frac{1}{1 - \alpha}$（資結書第~12.4~節）；
同樣的 $\alpha$ 下，探測次數比線性探測少很多，表可以裝得比較滿。 代價是每個鍵要多算一個雜湊
函數，而且探測位置跳來跳去，快取效果比線性探測差。

\textbf{(b) 不互質的問題。} 設 $d = \gcd(m, h_2(k)) > 1$。 $h(k, i) = h(k, j)$ 等價於 $m$ 整除
$(j - i)\,h_2(k)$，也就是 $\frac{m}{d}$ 整除 $j - i$。 所以探測序列走 $m/d$ 步就回到起點，之後
不斷重複，只會碰到 $m/d$ 個格子，其餘格子永遠不會被試到。 後果：
\begin{itemize}
\item 插入時，這 $m/d$ 格都滿了就找不到空位，即使表裡其他地方還是空的；程式若只檢查「試了
$m$ 次」才放棄，會在同幾格之間白白繞圈。
\item 搜尋不存在的鍵時，也只檢查了一部分格子，迴圈若以「回到起點」或「試滿 $m$ 次」結束，
前者過早結束、後者重複做白工。
\end{itemize}
例：表的大小 $m = 10$，某個鍵的 $h_1(k) = 3$、$h_2(k) = 4$，$\gcd(10, 4) = 2$。 探測序列是
$3$, $7$, $1$, $5$, $9$, $3$, $7$, $\ldots$，只會試 $5$ 個奇數格，偶數格全空也放不進去。

\textbf{(c) 解法。} 讓每個鍵的步長都和 $m$ 互質。 互質時探測序列的前 $m$ 項兩兩不同，正好走
遍全部 $m$ 格（資結書定理~12.4：若 $h(k, i) = h(k, j)$，$0 \le i < j < m$，則 $m$ 整除
$(j - i)\,h_2(k)$，互質推得 $m$ 整除 $j - i$，矛盾）。 常用的兩種取法：
\begin{itemize}
\item \textbf{$m$ 取質數}，$h_2(k)$ 的值落在 $1$ 到 $m - 1$ 之間，例如
$h_2(k) = 1 + (k \bmod (m - 1))$，或取一個比 $m$ 小的質數 $R$，$h_2(k) = R - (k \bmod R)$。
質數和 $1..m-1$ 的每個數都互質。
\item \textbf{$m$ 取 $2$ 的冪}，$h_2(k)$ 一律取奇數，例如 $h_2(k) = 2\,(k \bmod (m/2)) + 1$。
奇數和 $2^p$ 互質。
\end{itemize}
無論哪一種，$h_2(k)$ 都不能是 $0$（否則原地不動）。 上面的例子改成 $m = 11$、$h_2(k) = 4$，
序列是 $3, 7, 0, 4, 8, 1, 5, 9, 2, 6, 10$，$11$ 格都會試到。
\end{pastsoln}
